% Section owned by D (Thumula): strengths, assumptions, limitations, extensions.
% Drafted in full during the 21 Nov integration pass (B) from C's review-approved caveats
% (docs/c_paper_evidence.md) and our merged results/*.csv; D to review before the freeze.
\section{Critical discussion}\label{sec:critique}

\paragraph{The guarantees are conditional on the graph, and one number builds the graph.}
Everything tidy about the method --- real spectrum in $[-1,1]$, geometric decay,
a unique minimiser --- follows from the symmetric normalisation of $W$, not from
the labels. So the substantive modelling decision is $W$ itself: one RBF kernel
with a zero diagonal, whose single parameter $\sigma$ decides which edges exist
at all. Local and global consistency are assumptions \emph{about} that graph,
not properties readable off a dataset: connect two clusters of different class
and propagation will leak one into the other. Our sweep shows how sharply it
bites --- with $\alpha$, the data and the revealed labels fixed, mean LGC error
over the ten two-moons trials is $0.15\%$, $0.10\%$, $0.45\%$, $14.07\%$ and
$32.59\%$ at $\sigma=0.05,0.10,0.15,0.30,0.60$ (Table~\ref{tab:sweep}) --- and
the means hide the sharper point: at $\sigma=0.60$ the paired trials favour 1-NN
in eight of ten ($+10.2$ points against us), while at $\sigma\le0.30$ LGC wins
all ten. The defensible reading of the paper's headline is therefore not ``one
label per class suffices'' but ``one label per class suffices \emph{given a
neighbourhood scale that matches the geometry}''.

\paragraph{Selection is circular exactly where the method is meant to help.}
The authors acknowledge the difficulty of choosing parameters under extreme label
scarcity and report each algorithm at its own optimum, apart from $\alpha$
(Section~4). We did the only honest alternative for a toy study: fixed
$\sigma=0.15$ before the run, so nothing was tuned on evaluation error --- clean,
but arbitrary, and Table~\ref{tab:sweep} prices the arbitrariness. With two
revealed labels there is almost nothing to validate against, which makes
bandwidth stability the most useful follow-up to this paper.

\paragraph{$\alpha$ costs iterations, not accuracy.}
The contraction factor is $\alpha$, so hitting residual tolerance
$\mathrm{tol}$ costs error $\approx\mathrm{tol}/(1-\alpha)$ and roughly
$\ln(\mathrm{tol})/\ln(\alpha)$ iterations --- 1{,}405 measured at $\alpha=0.99$,
10{,}860 at $0.999$ (Table~\ref{tab:convergence}, where the bounds hold but are
loose). Toy accuracy rose monotonically over the same range
($0.77\to1.00$): $\alpha$ behaves as the advertised local-versus-global knob, and
on a real graph it decides affordability.

\paragraph{The linear algebra is dense.}
Building $W$ costs $O(n^{2})$ time and memory and the closed form needs a dense
solve: at $n=10^{5}$ the affinity alone is $10^{10}$ entries, roughly 80\,GB in
float64. Past the paper's own scale, sparse $k$-NN graphs with an iterative
solve, or its Figure~4 out-of-sample route, are the ways forward --- neither of
which we implemented.

\paragraph{One printed identity is wrong, harmlessly.}
Our audit (Appendix~\ref{app:variants}) finds Section~2's transpose claim false
as matrix algebra (relative gap $0.82$) while the corrected form, carrying an
extra $D$, verifies to $1.4\times10^{-15}$: a positive per-row rescaling leaves
the argmax untouched, and no variant moved more than $6/200$ toy predictions ---
a robustness finding as well as a correction.

\paragraph{What our evidence is not.}
Ours is 400 synthetic points, one revealed label per class, scored transductively
on unrevealed points only: a qualitative check of the mechanism, not of the digit
and text studies, and not a test on vectors unseen during graph construction.
Demanding a label from every class is favourable --- natural arrival does not
guarantee it --- and our ten seeds vary data and label placement together,
confounding those sources of variation. The spread we quote is the sample SD
across seeds, not a confidence interval; the paper's Figure~5 has no error bars
at all, which limits how much its small gaps between curves mean.
